\section{Conclusions and the next experiment}
\label{sec:conclusions}
\subsection*{What the evidence establishes}
The initial hypothesis that larger models would predict better is not supported within two A10 hours per run. Small wins the equal-dropout baseline and the final comparison of validation-selected settings. This is a useful negative result about resource allocation: more capacity is not sufficient to improve this pipeline. It is not a general claim that small models outperform large models. Size, updates, repeated exposure, and the suitability of the shared learning rate all change together.

The follow-up gives a more specific result. Raising dropout improves medium/large best validation losses by 0.0146/0.0154, while reducing their final losses by 0.0840/0.0919. Its main effect is therefore to reduce late worsening, with smaller gains at the checkpoints actually selected. Dropout does not close the baseline disadvantage. For small, reducing the train--validation gap produces no improvement in best validation loss. The intervention's value depends on the measured prediction objective, not on making curves look closer. One seed cannot tell us whether small's 0.0013 change is reproducible.

\subsection*{Why more training is not automatically the answer}
The baseline processes roughly 125.4/40.1/15.2 training-corpus equivalents for small/medium/large. These are exposure ratios under replacement sampling, not complete ordered epochs or new musical information. Muennighoff et al. find diminishing returns from repeated data in data-constrained language models \citep{muennighoff2023data}. This makes diversity a plausible limitation to investigate here. Their text-model repetition thresholds do not transfer directly to our event vocabulary, overlapping windows, or small architectures. Our curves do not prove a data-scaling law, and the larger models' lower training losses do not by themselves distinguish memorization from optimization effects.

Likewise, the shared time-based schedule equalizes the learning rate at the same elapsed time, not after the same number of updates. Validation and saving advance that clock, and ten-minute checks may miss each model's true best state. Reused validation windows can also favor their particular contents. Additional seeds, broader validation coverage, and comparable per-size tuning budgets are needed before treating the ranking as stable. The 177-piece consistency check addresses evaluation composition, not these training uncertainties.

\subsection*{Which future iteration is best justified?}
\textbf{First, improve data use on the small baseline.} Test training-only pitch transposition while leaving validation/test pieces unchanged. Shift note starts and matching endings together; reject shifts outside the playable range rather than clipping pitches. This preserves intervals while varying absolute pitch, but assumes transposed performances remain useful examples and may alter key/register statistics. Huang et al. use pitch/time augmentation in their piano experiments \citep[Section~4.2]{huang2019musictransformer}; this supports trying it, not a predicted gain for our tokenizer. Test augmentation separately from architecture changes so any improvement remains interpretable.

\textbf{For an architecture change, prefer a small relative-attention model over MoE or retrieval.} Hoffmann et al.'s improved-architecture discussion covers conditional computation and retrieval, rather than establishing a winner for limited music data \citep[Section~2]{hoffmann2022compute}. MoE reduces the active fraction of weights but adds routing and expert-management costs; Switch Transformers documents complexity and training-instability concerns \citep{fedus2022switch}. Our larger dense models already fit training well, so adding expert capacity does not directly address the observed failure to generalize. Retrieval would also require a useful music index and a training-only retrieval boundary; indexing the same small corpus does not create new performances. Neither option is ruled out, but both add less-targeted complexity at this stage.

Relative attention instead lets the small model use distances between events rather than relying only on absolute token positions. Music Transformer reports benefits on expressive piano event sequences \citep{huang2019musictransformer}, a closer task match than the large text-model evidence for MoE/retrieval. This is our preferred architectural hypothesis, not a demonstrated remedy for overfitting. Event distance is not elapsed musical time, and altered attention could lose optimized-kernel speed or use more memory. Start with the same 1,024-token context; do not simultaneously increase context or change dropout.

\textbf{Make the next comparison falsifiable.} First repeat the small baseline, then compare augmentation alone and relative attention alone, each with the same two-hour A10 budget, several seeds, identical validation windows, and comparable tuning effort. Report best validation loss, throughput, memory, parameter count, and actual tokens processed. A loss gain that disappears under equal time because attention is slower would reject its practical benefit for this project. Follow with a separately scheduled equal-token comparison to study learning efficiency. Check note/pitch families, tails/EOS, and matched-prompt MIDI continuations as secondary diagnostics. Because test outputs have been inspected, further choices must use validation; fresh confirmation needs a new holdout. These are proposed experiments, not completed results.
